\section{What repeated Newton queries can reuse}\label{sec:reuse}
The preceding direct sums require independent information in the specified
output.  A large number of central-path movements alone is insufficient.
We first give an explicit sparse KKT counterexample and then quantify one
way in which low-dimensional solution structure permits reuse.  Neither
statement treats residual certification, dense output, or state preparation
as free computation.

\subsection{A fixed sparse KKT state at every central multiplier}
Let $M\in\R^{n\times n}$ be invertible, with bidirectional sparse access,
row and column sparsity at most $s$, and $\norm M\leq1$.
For the one-cone program \eqref{eq:one-cone-program}, write
$v=M^{-1}e$, $a=\rho(\eta)$, and $q_a=1-a^2$.
Retain both $u=Mx$ and the equality multiplier $\lambda$.

\begin{proposition}[A common KKT ray]\label{prop:fixed-kkt}
At every exact central point, the augmented predictor equations are
\begin{equation}\label{eq:fixed-kkt}
 \begin{pmatrix}
 D_a&0&I\\0&0&-M^T\\I&-M&0
 \end{pmatrix}
 \begin{pmatrix}\dot u\\\dot x\\\dot\lambda\end{pmatrix}
 =\begin{pmatrix}2e\\0\\0\end{pmatrix},
 \qquad D_a=\frac2{q_a}I+\frac{4a^2}{q_a^2}ee^T.
\end{equation}
Their unique solution and normalized state are
\begin{equation}\label{eq:fixed-ray}
 (\dot u,\dot x,\dot\lambda)
 =\frac{q_a^2}{1+a^2}(e,v,0),\qquad
 \frac{(e,v,0)}{\sqrt{1+\norm v^2}}.
\end{equation}
The KKT matrix is at most $(s+2)$-sparse, with constant-overhead sparse-access
simulation from $M$ and $M^T$.  Thus any one-shot lower bound for preparing
$|v\rangle$ at sufficiently small constant state error transfers to each
individual KKT checkpoint, up to constants.
\end{proposition}
\begin{proof}
Stationarity in $x$ gives $M^T\lambda=0$, so $\lambda=0$.
The other stationarity equation is $u/(1-\norm u^2)=\eta e$;
hence $u=ae$ and $x=av$.  Differentiating gives
$a'=q_a^2/(1+a^2)$ and \eqref{eq:fixed-kkt}.
The second block row forces $\dot\lambda=0$, the third forces
$\dot u=M\dot x$, and $D_ae=2(1+a^2)e/q_a^2$ proves
\eqref{eq:fixed-ray}.  The remaining entries are public diagonal or
identity entries because $e$ is a basis vector.  This proves sparsity and
oracle simulation.  Since $\norm v\geq1$, the $x$ register has probability
at least $1/2$; measuring that register extracts $|v\rangle$ with constant
probability.  Sufficiently small constant Euclidean state error perturbs
this probability and the conditional state by only constant amounts.
\end{proof}

This is a state-output comparison, not a claim that every scalar derived
from that state is hard.  It also imposes no uniform bound on the condition
number of the augmented KKT matrix near the boundary.

The converse output distinction can be equally sharp.  A signed equality
chain $r_0=\alpha$, $r_i=\sigma_i r_{i-1}$, with $|\alpha|\leq1$
and objective $\alpha+2r_L$, reduces to maximizing
$(1+2H)\alpha$, where $H=\prod_i\sigma_i$.
Its optimum is three for $H=1$ and one for $H=-1$; two-sided relative
value error below $1/2$ therefore costs $\Theta(L)$ quantum or classical
queries.  Yet its nonzero reduced-coordinate solution state lives in a
one-dimensional Hilbert space and is input independent up to global phase.
The discarded sign or scale is essential to scalar readout.  The terminal
optimizer coordinate $r_L^*$ is actually one in both cases, so that
coordinate cannot serve as the parity readout.  The paired interval barrier
has parameter one, but this does not supply the hidden reduced objective
coefficient without acquiring the chain signs.

The same fixed instance can require arbitrarily many bounded local
movements.  For $\phi(u)=-\log(1-\norm u^2)$, direct inversion of its
radial Hessian gives
\[
 \norm{\nabla\phi(u)}_{u,*}^2
   =\frac{2\norm u^2}{1+\norm u^2}\leq1.
\]
Objective gap at most $0<\epsilon\leq1$ implies
$e^Tu\geq1-\epsilon/2$, so $1-\norm u^2\leq\epsilon$ and
$\phi(u)-\phi(0)\geq\log(1/\epsilon)$.
The gradient bound makes the barrier-metric length of every interior path
from the center to such a point at least $\log(1/\epsilon)$.
A chord with starting Dikin norm at most $R<1$ has length at most
$\log(1/(1-R))$.  One can see this directly for the ball: along a line
$\phi(u+td)$ is the sum of two negative logarithms of affine factors.
If $h^2=d^T\nabla^2\phi(u)d$, the absolute logarithmic slopes at zero
have squared sum $h^2$; at $t\geq0$ their denominators are at least
$1-th$.  Therefore the metric speed is at most $h/(1-th)$, whose
integral from zero to one is $-\log(1-h)$.
A round containing at most $m$ such chords consequently obeys
\begin{equation}\label{eq:movement-reuse}
 \#\text{rounds}\ \geq
 \frac{\log(1/\epsilon)}{m\log(1/(1-R))}.
\end{equation}
Pullback by $M$ preserves this barrier metric on the equality slice.
Equation~\eqref{eq:movement-reuse} concerns this explicit primal barrier
and bounded-Dikin chords; it is not an arbitrary algorithmic iteration
lower bound.

\begin{lemma}[An elementary hard ray]\label{lem:tree-hard-ray}
There are four-sparse matrices $M_\sigma$, of dimension $\Theta(L)$ and
condition number $\Theta(L)$, for which preparing
$|M_\sigma^{-1}e\rangle$ to sufficiently small constant state error has
quantum query complexity $\Theta(L)$.  One classical read-and-solve step
also uses $O(L)$ queries and arithmetic.  These matrices may be used in
Proposition~\ref{prop:fixed-kkt}.
\end{lemma}
\begin{proof}
Take a path $p_0,\ldots,p_L$ with hidden edge signs
$\sigma_1,\ldots,\sigma_L$.  Attach a complete public binary tree with
$h=2^{\lceil\log_2L\rceil}$ leaves at $p_0$, and a disjoint copy at
$p_L$.  For every nonroot node $z$ impose
$y_z-s_zy_{\pi(z)}=0$, where $s_z$ is the appropriate hidden path
sign or the public value one; impose $y_{p_0}=1$.
These equations define a triangular matrix $A_\sigma$ with row sparsity
two and column sparsity four.  Its dimension is
$N_L=L+4h-3$, and $\norm A\leq\sqrt{\norm A_1\norm A_\infty}
\leq\sqrt8$.  Set $M=A/\sqrt8$.
The unscaled solution $A^{-1}e$ is one on the first tree and
$H=\prod_i\sigma_i$ on the second, with prefix parities on the path.
The solution $M^{-1}e=\sqrt8 A^{-1}e$ has the same normalized state,
since its common public factor cancels. This state has constant probability
$2h/N_L\geq2/5$ in
the span of the two public uniform leaf states.  Measuring their sum and
difference identifies $H$ whenever that subspace is observed.
A constant number of sufficiently accurate preparations therefore computes
parity, which needs $\Omega(L)$ queries \citep{BealsEtAl2001}.
Reading all signs and forward substitution gives the matching upper bound.

Every entry of $A^{-1}$ is zero or a signed ancestor indicator. Its absolute row
and column sums are $O(L)$, so $\norm{A^{-1}}=O(L)$.
The path submatrix of $A^{-1}$ is, after diagonal sign changes, the
cumulative-sum matrix; applying it to the uniform vector shows norm
$\Omega(L)$.  Since $1\leq\norm A=O(1)$, the condition is
$\Theta(L)$.  The augmented KKT graph also has a width-three tree
decomposition: for each nonroot $z$ use the bag
$\{\lambda_z,u_z,x_z,x_{\pi(z)}\}$ and join bags along the dependency
tree; omit the absent parent variable at the root.
\end{proof}

After one $O(L)$-query solve, the stored vector $v$ determines every center
and predictor in \eqref{eq:fixed-ray} by public scalar rescaling, regardless
of the movement count in \eqref{eq:movement-reuse}.  Implicit classical
output then needs only a scalar and the stored vector.  Repeated dense
writes or quantum state preparations still cost work, but no further raw
input query is required.  More generally, reading all $O(ns)$ addressable
nonzero coefficient words of a fixed sparse matrix gives a finite cap on
subsequent coefficient-query complexity.  This cap applies to outputs
mathematically determined by those words, with arbitrary storage and
computation.  It does not reconstruct unspecified behavior of an arbitrary
unitary completion, nor bound arithmetic, gate, memory, or bit complexity.

The same reuse is directly relevant to scalar output.  For any fixed public
readout $c$, one has $c^Tx(\eta)=\rho(\eta)c^TM^{-1}e$.
One additive-$\epsilon$ estimate of $c^TM^{-1}e$ therefore supplies the
entire sequence of these scalar centers to additive error $\epsilon$,
since $0\leq\rho(\eta)<1$.  Choosing the plateau readout of
Section~\ref{sec:realizations} makes this one scalar classically hard;
its reuse needs no transfer from a state-output lower bound.

The sharper one-shot state lower
$\Omega(\kappa\log(1/\varepsilon_{\rm state}))$ of
\citet[Theorem~1]{MoriEtAl2026} holds for a constant-sparsity family under
coherent sparse position/value access, with
$0<\varepsilon_{\rm state}\leq1/11$ measured in Euclidean distance
between normalized solution states; right-hand-side preparation cost is
treated as negligible. This concerns a different output from our scalar
inverse forms. Its existence cannot justify
multiplying one-shot hardness by a movement count.  The elementary family
above already disproves that general inference without relying on a
precision-sensitive state-preparation theorem.

\subsection{Residual-certified span reuse}
Consider any sequence of nonsingular systems $H_tz_t=f_t\ne0$ and unit
rays $\psi_t=z_t/\norm{z_t}$.  Maintain a span of stored checkpoint
vectors.  At time $t$, minimize the \emph{current} relative residual
$\norm{H_ty-f_t}/\norm{f_t}$ over this span.  If the residual is too
large, solve the system once and append its ray.  Every use of an old
vector requires its image under the current operator; it is not assumed
that old residuals remain valid.

\begin{proposition}[Exact span]\label{prop:exact-reuse}
If the full sequence of solution rays has span dimension $r$, exact
residual minimization and exact checkpoint rays require at most $r$
refresh solves, for any acceptance threshold in $(0,1)$.  This is sharp.
\end{proposition}
\begin{proof}
If the current ray belongs to the stored span, that span contains the exact
solution and its minimum residual is zero.  Each rejection therefore adds
a linearly independent ray.  For sharpness take $H_t=I$, $f_t=e_t$,
$t=1,\ldots,r$; every not-yet-stored vector has relative residual one.
\end{proof}

For a robust version, set
$A_t=(\norm{z_t}/\norm{f_t})H_t$, so
$A_t\psi_t=f_t/\norm{f_t}$, and assume $\norm{A_t}\leq B_0$.
Necessarily $B_0\geq1$.  These normalized operators are only for
analysis; the algorithm can minimize the original relative residual and
need not know $\norm{z_t}$.  Let $0<\eta\leq1$ be the desired residual
scale.  Suppose minimization returns a candidate with residual at most
$\eta/8$ above the true minimum, and the test estimates that candidate's
residual to additive error $\eta/8$.  Accept if the reported value is at
most $3\eta/4$.  An accepted candidate has true relative residual at
most $7\eta/8$, whereas rejection implies minimum residual greater than
$\eta/2$.

\begin{theorem}[Numerical span with normalized checkpoints]\label{thm:robust-reuse}
Let $D$ be the system dimension and let $1\leq r\leq D$ be an integer.
Suppose every refreshed stored vector $\widehat\psi_t$ has norm one and
$\norm{\widehat\psi_t-\psi_t}\leq\xi$.  Define the Kolmogorov width
\[
 d_r=\inf_{\dim S=r}\sup_t\operatorname{dist}(\psi_t,S),
 \qquad \Delta=\frac{\eta}{2B_0}-\xi>0.
\]
If
\begin{equation}\label{eq:width-refresh}
 d_r+\xi<\frac{\Delta^{2-1/r}}{2r},
\end{equation}
then the residual procedure makes at most $2r-1$ refreshes, on the event
that all stated minimization, test and checkpoint errors hold.
\end{theorem}
\begin{proof}
For a stored span $V$, its best relative residual is at most
$B_0\operatorname{dist}(\psi_t,V)$.  Rejection therefore implies
$\operatorname{dist}(\psi_t,V)>\eta/(2B_0)$, and the next approximate
stored vector has distance greater than $\Delta$ from $V$.
Suppose $2r$ refreshes occur and form a matrix $Y$ from their normalized
columns.  Its first QR diagonal has magnitude one and all later diagonals
exceed $\Delta$, giving
\[
 \sqrt{\det(Y^*Y)}>\Delta^{2r-1}.
\]
Choose an $r$-space $S$ with all column distances at most
$d_r+\xi+\epsilon$, where $\epsilon>0$ is arbitrary.
The first $r$ singular values of $Y$ are at most $\sqrt{2r}$ since its
columns have norm one.  The remaining $r$ are at most
$\norm{(I-\Pi_S)Y}\leq\sqrt{2r}(d_r+\xi+\epsilon)$, since
$\Pi_SY$ has rank at most $r$.  Thus
\[
 \sqrt{\det(Y^*Y)}\leq(2r)^r(d_r+\xi+\epsilon)^r.
\]
Letting $\epsilon\downarrow0$ contradicts \eqref{eq:width-refresh}.
Normalization of every stored column is used in this volume bound.
\end{proof}

There is a stronger bound under a correspondingly stronger stability
assumption.  Suppose the exact rays have rank $r$, and whenever a new exact
ray is dependent on \emph{the exact versions of the previously selected
checkpoint rays}, it has a representation in those rays whose coefficients
have sum of absolute values at most $\Lambda$.  If
$B_0\Lambda\xi\leq\eta/2$, replacing that representation by the stored
approximate rays has relative residual at most $\eta/2$.
It cannot trigger a rejection.  Hence at most $r$ refreshes occur.
Stability in an unspecified final basis is insufficient for this adaptive
argument.  This condition permits checkpoint error linear in $\eta$;
any resulting tomography complexity must still use the cost law of the
particular tomography interface.

These statements specialize the established ideas of subspace recycling
and greedy reduced bases \citep{ParksEtAl2006,BinevEtAl2011} to a
residual-certified checkpoint contract.  They give refresh counts, not
unconditional quantum speedups.  With explicit vectors of dimension $D$
and $R$ retained checkpoints, forming current images and solving their
least-squares problem costs, for example,
$O(R\,\mathrm{mv}(H_t)+DR^2+R^3)$ arithmetic per test.
Quantum tests have their own input preparation, normalization and
coefficient-conditioning costs.  If the systems are expressed in a linear
nullspace representation enforcing homogeneous feasibility equations,
span combinations preserve those equations.  For an affine feasibility
constraint one must instead retain a particular feasible point and combine
homogeneous directions.  Neither observation supplies cone-scaling access
for free.

\subsection{When low numerical rank is justified, and what it does not buy}
A useful sufficient regularity condition is holomorphic dependence.  Suppose
a unit-ray curve on $[-1,1]$ extends holomorphically to a Bernstein ellipse
of parameter $\chi>1$, continuously to its boundary, with norm at most
$W$.  Its vector-valued Chebyshev coefficients satisfy
$\norm{c_j}\leq2W\chi^{-j}$ for $j\geq1$, by the usual contour
integral on the ellipse.  The degree-$(r-1)$ truncation lies in the span
of $r$ coefficient vectors, and its tail is bounded by
\begin{equation}\label{eq:analytic-width}
 d_r\leq\sum_{j=r}^{\infty}2W\chi^{-j}
       =\frac{2W\chi^{-(r-1)}}{\chi-1}.
\end{equation}
Thus a uniform ellipse and bound give a logarithmic numerical-rank scale.
They are actual hypotheses, not consequences of strict complementarity.
For example, the strictly complementary separable QP
$\min_{x\geq0}\frac12\norm x^2+c^Tx$, with distinct $c_i>0$, has
barrier central coordinates and predictors
\[
 x_i(\tau)=\frac{-c_i+\sqrt{c_i^2+4\tau}}2,
 \qquad x_i'(\tau)=\frac1{\sqrt{c_i^2+4\tau}}.
\]
The coordinate functions of the predictor are linearly independent:
analytic continuation of a vanishing linear combination to a neighborhood
of $-c_i^2/4$ forces its $i$th coefficient to vanish, since only that
summand has that branch singularity.  The normalized predictor rays
therefore span the full dimension; common normalization does not change
linear relations among coordinate functions on the real interval.
Nor is there a uniform complex neighborhood as $\min_i c_i\downarrow0$.
For a normalized-ray example in dimension two, choose $c_2=2c_1$.
The squared first coordinate of the normalized predictor is
$(c_2^2+4\tau)/(c_1^2+c_2^2+8\tau)$, with a nonremovable pole tending
to zero from below.  A uniform Bernstein ellipse for normalized rays on
$[0,1]$ is therefore impossible in this family.

Low rank alone also does not make certification cheap.  With coefficient
queries, take $H_0=H_1=I_D$, $f_0=e_0$, and
\[
 f_1=e_0+\sum_{i=1}^{D-1}a_i e_i,\qquad a_i\in\{0,1\},
 \quad\sum_i a_i\in\{0,1\}.
\]
The solution span has dimension at most two and $B_0=1$.
The minimum relative residual from $\operatorname{span}\{e_0\}$ is zero
or $1/\sqrt2$.  Any test separating these cases solves unique search,
requiring $\Omega(D)$ randomized or $\Omega(\sqrt D)$ quantum queries.
For completeness, under a uniform marked location a $Q$-query randomized
algorithm hits it with probability at most $Q/(D-1)$ until the transcript
changes; the quantum hybrid bound changes the average final-state distance
by at most $O(Q/\sqrt D)$.
An exact norm oracle for $f_1$ would reveal the case immediately, which is
why the coefficient-query input is explicit here.

Even a single stored ray does not eliminate explicit-output cost.  For
$H=I$ and $f_b=D^{-1/2}((-1)^{b_1},\ldots,(-1)^{b_D})$, a phase query
implemented by the canonical bit oracle prepares the amplitude state.
The input for the following lower bound is this bit/value oracle; a
phase-only unitary would not expose the overall sign.
A classical vector within Euclidean error
$1/4$ determines all but at most $D/16$ signs by rounding its coordinates.
This requires $\Omega(D)$ classical queries and $\Omega(D)$ writes;
its quantum query cost is also $\Omega(D)$.  To see the latter without
assuming exact recovery, a $Q$-query final state, as the signs vary, lies
in a space of dimension at most $\sum_{j=0}^Q\binom Dj$ by the
multilinear query expansion \citep{BealsEtAl2001}.
Any one rounded output is correct for at most
$V=\sum_{j\leq D/16}\binom Dj\leq2^{H_2(1/16)D}$ inputs.
For any measurement, its average success on uniform inputs is at most
$V\sum_{j\leq Q}\binom Dj/2^D$: bound each state's density matrix by
the projector onto this common space and sum the measurement operators.
Constant success consequently requires $Q=\Omega(D)$.
Here $H_2$ is binary entropy.  The argument separates inexpensive state
preparation from a classical vector, even when the same rank-one target
is requested at every checkpoint.
